\documentclass[11pt,a4paper]{article}
% >>> PERSONALIZZA <<<
\newcommand{\titolobreve}{Basi di dati \,·\, Il modello Entity-Relationship}  % testo dell'intestazione
\newcommand{\etichettasezione}{Capitolo}                  % oppure: Sezione
\newcommand{\titolopdf}{Il modello Entity-Relationship – Manuale di studio}       % metadati PDF + copertina
\newcommand{\autore}{Dalle lezioni di Basi di dati (P. Atzeni), A.A. 2026–2027}                         % metadati PDF + copertina

% ---------------- Matematica (Palatino, stile manuale scientifico) ----------------
\usepackage{amsmath,amssymb,mathtools}
\usepackage[sc]{mathpazo}
\usepackage{bm}
\linespread{1.06}

% ---------------- Font del testo ----------------
\usepackage[no-math]{fontspec}
\setmainfont{TeX Gyre Pagella}
\setsansfont{Fira Sans}[BoldFont={Fira Sans SemiBold},ItalicFont={Fira Sans Italic},BoldItalicFont={Fira Sans SemiBold Italic}]
\newfontfamily\headfont{Fira Sans}[BoldFont={Fira Sans SemiBold},ItalicFont={Fira Sans Italic}]
\newfontfamily\headlight{Fira Sans Light}
\newfontfamily\headblack{Fira Sans ExtraBold}

\usepackage{babel}
\babelprovide[import,main]{italian}

% ---------------- Pagina ----------------
\usepackage[a4paper,top=2.5cm,bottom=2.4cm,left=2.3cm,right=2.3cm,headheight=15pt,headsep=0.9cm]{geometry}
\usepackage[table]{xcolor}
\usepackage{graphicx,tikz,booktabs,array,tabularx,enumitem,ragged2e,pifont,microtype}
\usetikzlibrary{calc,positioning,shapes.geometric,arrows.meta,shadows}
\usepackage[most]{tcolorbox}
\usepackage{titlesec,fancyhdr,tocloft}
\usepackage[hidelinks]{hyperref}
\usepackage{needspace,etoolbox}
\AtBeginDocument{\hypersetup{pdftitle={\titolopdf},pdfauthor={\autore},pdfsubject={\titolobreve}}}

\setlength{\parindent}{0pt}
\setlength{\parskip}{5pt plus 1pt}
\raggedright
\setlength{\RaggedRightParindent}{0pt}
\emergencystretch=2em
\widowpenalty=10000 \clubpenalty=10000 \displaywidowpenalty=10000

% ---------------- Palette ----------------
\definecolor{navy}{HTML}{1F3A5F}
\definecolor{teal}{HTML}{1B8A8A}
\definecolor{blu}{HTML}{2A6FDB}
\definecolor{rosso}{HTML}{E0533D}
\definecolor{ambra}{HTML}{D9940F}
\definecolor{viola}{HTML}{6B4FBB}
\definecolor{verde}{HTML}{2E9E5B}
\definecolor{grigio}{HTML}{5A6170}
\definecolor{chiaro}{HTML}{EEF2F7}
\definecolor{testo}{HTML}{22262D}
\color{testo}

% vettori in grassetto
\newcommand{\x}{\mathbf{x}}
\newcommand{\vv}{\mathbf{v}}
\newcommand{\mmu}{\bm{\mu}}
\newcommand{\SSigma}{\bm{\Sigma}}
\newcommand{\argmax}{\operatorname*{arg\,max}}
\newcommand{\T}{^{\mathsf{t}}}
\newcommand{\hl}[1]{{\color{navy}\bfseries #1}}
\newcommand{\key}[1]{\textcolor{teal}{\textbf{#1}}}
\newcommand{\slide}[1]{{\headfont\footnotesize\color{grigio}\,[slide #1]}}

% ---------------- Liste ----------------
\setlist{itemsep=2pt,topsep=3pt,parsep=0pt}
\setlist[itemize,1]{label={\color{teal}\small$\blacktriangleright$},leftmargin=1.3em}
\setlist[itemize,2]{label={\color{ambra}\footnotesize$\bullet$},leftmargin=1.2em}
\setlist[enumerate,1]{label={\headfont\bfseries\color{navy}\arabic*.},leftmargin=1.6em}

% ---------------- Icone (tikz) ----------------
\newcommand{\iconbulb}{\tikz[baseline=-0.6ex,scale=.85]{\fill[white](0,.06)circle(.13);\fill[white](-.07,-.15)rectangle(.07,-.06);\draw[white,line width=.4pt](-.06,-.19)--(.06,-.19);}}
\newcommand{\iconwarn}{\tikz[baseline=-0.6ex]{\node[regular polygon,regular polygon sides=3,fill=white,inner sep=1.1pt,minimum size=11pt]{}; \node[font=\headblack\tiny,text=rosso] at (0,-.02){!};}}
\newcommand{\iconml}{\tikz[baseline=-0.6ex,scale=.8]{\foreach \y in {-.16,.16}{\fill[white](-.2,\y)circle(.05);} \foreach \y in {-.2,0,.2}{\fill[white](.2,\y)circle(.05);\draw[white,line width=.35pt](-.2,.16)--(.2,\y);\draw[white,line width=.35pt](-.2,-.16)--(.2,\y);}}}

% ---------------- Box ----------------
\tcbset{every box/.style={halign=flush left}}
\tcbset{
  base/.style={enhanced,unbreakable,arc=3pt,boxrule=0pt,left=11pt,right=11pt,top=9pt,bottom=8pt,
    before skip=12pt,after skip=12pt,fonttitle=\headfont\bfseries\small,
    attach boxed title to top left={xshift=10pt,yshift=-9pt},
    boxed title style={arc=2.5pt,boxrule=0pt,left=6pt,right=6pt,top=2pt,bottom=2pt},
    overlay unbroken and first={}},
}
\DeclareTColorBox{intuizione}{ O{Intuizione} O{} }{base,colback=teal!6,colbacktitle=teal,coltitle=white,
  borderline west={3pt}{0pt}{teal},title={\iconbulb\ \ #1},top=13pt,#2}
\DeclareTColorBox{tecnica}{ O{In termini tecnici} O{} }{base,colback=navy!5,colbacktitle=navy,coltitle=white,
  borderline west={3pt}{0pt}{navy},title={$\Sigma$\ \ #1},top=13pt,#2}
\DeclareTColorBox{esempio}{ O{Esempio} O{} }{base,colback=ambra!7,colbacktitle=ambra,coltitle=white,
  borderline west={3pt}{0pt}{ambra},title={\ding{46}\ \ #1},top=13pt,#2}
\DeclareTColorBox{attenzione}{ O{Attenzione} O{} }{base,colback=rosso!6,colbacktitle=rosso,coltitle=white,
  borderline west={3pt}{0pt}{rosso},title={\iconwarn\ \ #1},top=13pt,#2}
\DeclareTColorBox{mlbox}{ O{Approfondimento} O{} }{base,colback=viola!6,colbacktitle=viola,coltitle=white,
  borderline west={3pt}{0pt}{viola},title={\iconml\ \ #1},top=13pt,#2}
\newtcolorbox{formula}[1][]{enhanced,arc=5pt,boxrule=0.9pt,colframe=navy!55,colback=white,
  left=8pt,right=8pt,top=6pt,bottom=6pt,before skip=10pt,after skip=10pt,halign=flush center,
  #1}
\newtcolorbox{formulat}[1]{enhanced,arc=5pt,boxrule=0.9pt,colframe=navy!55,colback=white,
  left=8pt,right=8pt,top=10pt,bottom=6pt,before skip=14pt,after skip=10pt,halign=flush center,
  attach boxed title to top center={yshift=-8pt},boxed title style={colback=navy!55,arc=2pt,boxrule=0pt,top=1pt,bottom=1pt},
  fonttitle=\headfont\bfseries\footnotesize,coltitle=white,title={#1}}
\newtcolorbox{appunti}[1]{enhanced,unbreakable,arc=4pt,boxrule=.8pt,colframe=grigio!60,colback=chiaro,
  left=12pt,right=12pt,top=12pt,bottom=10pt,before skip=16pt,after skip=8pt,
  fonttitle=\headfont\bfseries\small,coltitle=white,colbacktitle=grigio,
  attach boxed title to top left={xshift=10pt,yshift=-9pt},boxed title style={arc=2.5pt,boxrule=0pt,left=6pt,right=6pt},
  title={\ding{43}\ \ #1},fontupper=\small}
\newcommand{\pag}[1]{\par\vspace{4pt}{\headfont\bfseries\scriptsize\setlength{\fboxsep}{1.5pt}\colorbox{teal}{\color{white}\,p.\,#1\,}}\enspace}

% ---------------- Titoli ----------------
\newcounter{sez}
\newcommand{\sezione}[3]{% num titolo slide
  \clearpage
  \refstepcounter{section}
  \addcontentsline{toc}{section}{\protect\numberline{\thesection}#2}
  \markboth{#2}{#2}
  \begin{tikzpicture}
    \fill[navy,rounded corners=8pt] (0,0) rectangle (\textwidth,2.6cm);
    \fill[teal,rounded corners=8pt] (0,0) rectangle (2.6cm,2.6cm);
    \fill[teal] (1.6cm,0) rectangle (2.6cm,2.6cm);
    \node[font=\headblack\fontsize{40}{40}\selectfont,text=white] at (1.3cm,1.3cm) {\thesection};
    \node[anchor=south west,font=\headfont\bfseries\fontsize{19}{22}\selectfont,text=white,text width=\textwidth-3.6cm,align=left] at (3.0cm,.85cm) {#2};
    \node[anchor=west,font=\headlight\small,text=white!80!teal] at (3.1cm,.48cm) {\etichettasezione\ #1 \quad·\quad #3};
  \end{tikzpicture}
  \par\vspace{6pt}
}
\titleformat{\subsection}{\headfont\bfseries\large\color{navy}}{\colorbox{teal!15}{\color{teal}\thesubsection}}{0.6em}{}
\titlespacing*{\subsection}{0pt}{16pt}{6pt}
% un titolo non resta mai da solo in fondo alla pagina
\pretocmd{\subsection}{\Needspace{9\baselineskip}}{}{}
\pretocmd{\subsubsection}{\Needspace{6\baselineskip}}{}{}
\titleformat{\subsubsection}{\headfont\bfseries\normalsize\color{teal}}{}{0em}{}
\titlespacing*{\subsubsection}{0pt}{10pt}{3pt}

% ---------------- Header / footer ----------------
\pagestyle{fancy}
\fancyhf{}
\renewcommand{\headrulewidth}{0pt}
\fancyhead[L]{\headfont\footnotesize\color{grigio}\titolobreve}
\fancyhead[R]{\headfont\footnotesize\color{teal}\leftmark}
\fancyfoot[C]{\tikz{\node[circle,fill=navy,text=white,font=\headfont\bfseries\footnotesize,minimum size=17pt,inner sep=0pt]{\thepage};}}
\fancypagestyle{plain}{\fancyhf{}\fancyfoot[C]{\tikz{\node[circle,fill=navy,text=white,font=\headfont\bfseries\footnotesize,minimum size=17pt,inner sep=0pt]{\thepage};}}}

% ---------------- Indice ----------------
\renewcommand{\cfttoctitlefont}{\headfont\bfseries\Huge\color{navy}}
\renewcommand{\cftsecfont}{\headfont\bfseries\color{navy}}
\renewcommand{\cftsecpagefont}{\headfont\bfseries\color{teal}}
\renewcommand{\cftsecleader}{\color{grigio!50}\cftdotfill{\cftdotsep}}
\renewcommand{\cftsubsecfont}{\headfont\small\color{grigio}}
\renewcommand{\cftsubsecpagefont}{\headfont\small\color{grigio}}
\setlength{\cftbeforesecskip}{7pt}
\setcounter{tocdepth}{1}

% tabelle
\newcommand{\thd}[1]{{\headfont\bfseries\small\color{white}#1}}
\renewcommand{\arraystretch}{1.25}
\renewcommand{\tabularxcolumn}[1]{m{#1}}
\newcolumntype{Y}{>{\raggedright\arraybackslash}X}
\newcommand{\fig}[2][\linewidth]{\begin{center}\includegraphics[width=#1]{figs/#2}\end{center}}
% figura + didascalia sempre sulla stessa pagina: \figura[larghezza]{file.pdf}{didascalia}
\newcommand{\figura}[3][\linewidth]{\par\noindent\begin{minipage}{\linewidth}\fig[#1]{#2}\didascalia{#3}\end{minipage}\par}
\newcommand{\didascalia}[1]{\vspace{-8pt}\begin{center}\begin{minipage}{.9\linewidth}\centering\headfont\footnotesize\color{grigio}#1\end{minipage}\end{center}\vspace{-2pt}}


% ---------------- Codice e terminali ----------------
\setmonofont{DejaVu Sans Mono}[Scale=0.84]
\usepackage{listings}
\definecolor{termbg}{HTML}{14202E}
\definecolor{termfg}{HTML}{E4E9F0}
\definecolor{termprompt}{HTML}{5FD3B4}
\definecolor{termhl}{HTML}{F5C04A}
\definecolor{termdim}{HTML}{8A97A8}
\lstdefinestyle{term}{basicstyle=\ttfamily\small\color{termfg},columns=fullflexible,keepspaces=true,
  breaklines=false,
  showstringspaces=false,upquote=true,
  moredelim=[is][\color{termprompt}\bfseries]{§}{§},
  moredelim=[is][\color{termhl}\bfseries]{¤}{¤},
  moredelim=[is][\color{termdim}\itshape]{°}{°},
  literate={à}{{\`a}}1 {è}{{\`e}}1 {é}{{\'e}}1 {ì}{{\`i}}1 {ò}{{\`o}}1 {ù}{{\`u}}1}
\lstdefinestyle{file}{basicstyle=\ttfamily\small\color{testo},columns=fullflexible,keepspaces=true,breaklines=false,
  upquote=true,showstringspaces=false,
  moredelim=[is][\color{teal}\bfseries]{§}{§},
  moredelim=[is][\color{rosso}\bfseries]{¤}{¤},
  moredelim=[is][\color{grigio}\itshape]{°}{°},
  literate={à}{{\`a}}1 {è}{{\`e}}1 {é}{{\'e}}1 {ì}{{\`i}}1 {ò}{{\`o}}1 {ù}{{\`u}}1}
\tcbuselibrary{listings}
% terminale: argomento = titolo della finestra
\newtcblisting{terminale}[1]{enhanced,listing only,breakable,colback=termbg,colframe=termbg,arc=5pt,boxrule=0pt,
  left=9pt,right=9pt,top=4pt,bottom=4pt,before skip=9pt,after skip=9pt,
  listing options={style=term},
  title={\tikz[baseline=-.6ex]{\foreach \c [count=\i] in {rosso,ambra,verde}{\fill[\c](\i*.32,0)circle(.09);}}\hfill{\ttfamily\footnotesize\color{termdim}\detokenize{#1}}\hfill\phantom{\tikz{\fill(0,0)circle(.09);\fill(.64,0)circle(.09);}}},
  colbacktitle=termbg!80!white,coltitle=white,toptitle=3pt,bottomtitle=3pt,titlerule=0pt}
% file di configurazione: argomento = nome del file
\newtcblisting{filecfg}[1]{enhanced,listing only,breakable,colback=white,colframe=navy!40,arc=4pt,boxrule=.7pt,
  left=9pt,right=9pt,top=4pt,bottom=4pt,before skip=9pt,after skip=9pt,listing options={style=file},
  attach boxed title to top left={xshift=8pt,yshift=-8pt},boxed title style={colback=navy!70,arc=2pt,boxrule=0pt},
  fonttitle=\ttfamily\footnotesize\bfseries,coltitle=white,title={\detokenize{#1}},top=10pt}
\newcommand{\cmd}[1]{\texttt{\color{navy}#1}}
\newcommand{\inl}[1]{{\setlength{\fboxsep}{1.2pt}\colorbox{chiaro}{\ttfamily\small\color{navy}#1}}}
% scheda comando
\newtcolorbox{scheda}[2][]{enhanced,breakable,arc=4pt,boxrule=0pt,colback=white,
  borderline={.8pt}{0pt}{teal!60},left=11pt,right=11pt,top=10pt,bottom=8pt,before skip=14pt,after skip=12pt,
  title={\ttfamily\bfseries\normalsize #2},fonttitle=\color{white},colbacktitle=teal,
  toptitle=4pt,bottomtitle=4pt,#1}
\newcommand{\campo}[1]{\par\vspace{3pt}{\headfont\bfseries\footnotesize\color{teal}\MakeUppercase{#1}}\par\vspace{1pt}}
\tikzset{
  dev/.style={draw=navy,fill=white,line width=1pt,rounded corners=3pt,minimum width=1.25cm,minimum height=.75cm,font=\headfont\bfseries\small,text=navy},
  host/.style={dev,fill=blu!8,draw=blu},
  rtr/.style={circle,draw=teal,fill=teal!10,line width=1.1pt,minimum size=1.05cm,font=\headfont\bfseries\small,text=teal!70!black},
  cd/.style={draw=ambra,fill=ambra!15,line width=1pt,rounded corners=7pt,minimum width=1.0cm,minimum height=.62cm,font=\headfont\bfseries,text=ambra!60!black},
  lnk/.style={line width=1.2pt,color=grigio!80},
  lab/.style={font=\ttfamily\scriptsize,text=testo,fill=white,inner sep=1.2pt,rounded corners=1pt},
  ipl/.style={font=\ttfamily\scriptsize\bfseries,text=rosso!85!black,fill=white,inner sep=1.2pt},
  macl/.style={font=\ttfamily\scriptsize,text=viola,fill=viola!8,inner sep=1.5pt,rounded corners=1pt},
  netl/.style={font=\ttfamily\scriptsize\bfseries,text=verde!70!black,fill=verde!8,draw=verde!50,rounded corners=2pt,inner sep=2pt},
}
\newcommand{\wsrow}[7]{\rowcolor{#7}#1 & #2 & #3 & #4 & #5 & #6\\}

\newcommand{\sint}[1]{{\ttfamily\small\color{navy}\raggedright #1\par}}


% ---------------- Codice sincronizzato con i file (manuali con codice) ----------------
% Regola: il codice e l'output NON si copiano a mano nel .tex: si includono dai file veri.
\definecolor{pykw}{HTML}{1F3A5F}
\definecolor{pystr}{HTML}{2E7D4F}
\definecolor{pycom}{HTML}{6B7280}
\definecolor{pybg}{HTML}{F7F9FC}
\lstdefinestyle{py}{language=Python,basicstyle=\ttfamily\small\color{testo},columns=fullflexible,keepspaces=true,
  breaklines=false,upquote=true,showstringspaces=false,
  keywordstyle=\color{pykw}\bfseries,stringstyle=\color{pystr},commentstyle=\color{pycom}\itshape,
  morekeywords={as,with,None,True,False,self},emphstyle=\color{teal}\bfseries,
  numbers=left,numberstyle=\tiny\color{grigio},numbersep=7pt,xleftmargin=12pt,
  literate={à}{{\`a}}1 {è}{{\`e}}1 {é}{{\'e}}1 {ì}{{\`i}}1 {ò}{{\`o}}1 {ù}{{\`u}}1}
% parole da evidenziare in teal nei listati (classi/funzioni chiave dell'argomento)
\newcommand{\paroleevidenziate}[1]{\lstset{emph={#1}}}
\tcbset{codicestyle/.style={enhanced,listing only,colback=pybg,colframe=teal!55,arc=4pt,boxrule=.7pt,
  left=4pt,right=8pt,top=10pt,bottom=4pt,before skip=12pt,after skip=10pt,
  attach boxed title to top left={xshift=8pt,yshift=-8pt},boxed title style={colback=teal,arc=2pt,boxrule=0pt},
  fonttitle=\ttfamily\footnotesize\bfseries,coltitle=white},
  outputstyle/.style={enhanced,listing only,colback=termbg,colframe=termbg,arc=5pt,boxrule=0pt,
  left=9pt,right=9pt,top=4pt,bottom=4pt,before skip=9pt,after skip=9pt,
  colbacktitle=termbg!80!white,coltitle=white,toptitle=3pt,bottomtitle=3pt,titlerule=0pt,
  fonttitle=\ttfamily\footnotesize}}
% \codicefile[opzioni]{file.py}{prima}{ultima}: righe reali del file, numerate come nel file.
%   Fino a ~25 righe il blocco resta unito; oltre, passa [breakable] oppure dividi in più blocchi.
\newcommand{\codicefile}[4][]{\tcbinputlisting{codicestyle,listing file=code/#2,
  listing options={style=py,firstline=#3,lastline=#4,firstnumber=#3},
  title={\ding{230}\ \detokenize{#2}\ \ (righe #3–#4)},#1}}
% \codicefileintero[opzioni]{file.py}: tutto il file
\newcommand{\codicefileintero}[2][]{\tcbinputlisting{codicestyle,listing file=code/#2,
  listing options={style=py},title={\ding{230}\ \detokenize{#2}},#1}}
% \outputfile[opzioni]{output/file.txt}{titolo}: output salvato dall'esecuzione reale
\newcommand{\outputfile}[3][]{\tcbinputlisting{outputstyle,listing file=#2,listing options={style=term},
  title={\tikz[baseline=-.6ex]{\foreach \c [count=\i] in {rosso,ambra,verde}{\fill[\c](\i*.32,0)circle(.09);}}\hfill{\ttfamily\footnotesize\color{termdim}\detokenize{#3}}\hfill\phantom{\tikz{\fill(0,0)circle(.09);\fill(.64,0)circle(.09);}}},#1}}
% codice scritto direttamente nel .tex (solo frammenti brevi che non esistono come file)
\newtcblisting{codice}[1]{codicestyle,listing options={style=py},title={\ding{230}\ \detokenize{#1}}}
% richiamo di righe nel testo: \rr{36} oppure \rr{36}[41]  ->  r. 36–41
\NewDocumentCommand{\rr}{m o}{{\headfont\bfseries\scriptsize\setlength{\fboxsep}{1.4pt}%
  \colorbox{teal!15}{\color{teal!70!black}r.\,#1\IfValueT{#2}{–#2}}}}

% ---------------- Esercizi e soluzioni ----------------
\DeclareTColorBox{esercizi}{ O{Esercizi} O{} }{base,colback=blu!5,colbacktitle=blu,coltitle=white,
  borderline west={3pt}{0pt}{blu},title={\ding{45}\ \ #1},top=13pt,#2}
% nelle soluzioni (appendice): \soluzione{1.2}{testo}
\newcommand{\soluzione}[2]{\par\vspace{4pt}{\headfont\bfseries\small\color{blu}Esercizio #1.}\enspace #2\par}

% ---------------- Elementi per manuali senza formule/codice ----------------
\DeclareTColorBox{consiglio}{ O{Consiglio} O{} }{base,colback=verde!6,colbacktitle=verde,coltitle=white,
  borderline west={3pt}{0pt}{verde},title={\ding{51}\ \ #1},top=13pt,#2}
% passi numerati con cerchi (procedure, istruzioni d'uso)
\newlist{passi}{enumerate}{1}
\setlist[passi]{label={\protect\tikz[baseline=(n.base)]\protect\node[circle,fill=teal,text=white,font=\headfont\bfseries\footnotesize,inner sep=0pt,minimum size=1.35em](n){\arabic*};},leftmargin=2.1em,itemsep=6pt,topsep=6pt}
% box «In breve» per riassunti di poche righe in apertura di capitolo
\newtcolorbox{inbreve}{enhanced,arc=4pt,boxrule=0pt,colback=chiaro,left=12pt,right=12pt,top=8pt,bottom=8pt,
  borderline north={2pt}{0pt}{teal},fontupper=\small,before skip=6pt,after skip=10pt}
